\documentclass{article}
\usepackage[utf8]{inputenc}
\usepackage[russian]{babel}
\usepackage{amsmath}
\usepackage{amssymb}
\usepackage{graphicx}
\usepackage{hyperref}

\title{Отчет о проделанной работе: Методы Knowledge Graph Completion (KGC) для автоматического заполнения недостающих связей во временных графах знаний}
\author{}
\date{}

\begin{document}

\maketitle

### 2.3.1.4 Методы Knowledge Graph Completion (KGC) для автоматического заполнения недостающих связей в графе знаний

В данном разделе представлен обзор методов Knowledge Graph Completion (KGC), используемых в проекте для автоматического заполнения недостающих связей во временных графах знаний. Основное внимание уделяется архитектуре моделей, их рабочему процессу, используемому набору данных и результатам оценки.

#### 2.3.1.4.1 Архитектура модели и используемые методы

Проект базируется на двух ключевых компонентах: модели для дополнения временных графов знаний (Temporal Knowledge Base Completion, TKBC) и модели для ответов на вопросы (Question Answering, QA).

##### 2.3.1.4.1.1 Модель Temporal Knowledge Base Completion (TComplEx)

Для обучения эмбеддингов сущностей, отношений и временных меток используется модель TComplEx. TComplEx является тензорно-декомпозиционным подходом, который позволяет захватывать временную динамику графа знаний. Эмбеддинги, полученные с помощью TComplEx, представляют собой комплексные векторы, что позволяет модели эффективно моделировать различные типы связей и их эволюцию во времени. Процесс обучения TComplEx включает генерацию негативных примеров, получение комплексных эмбеддингов, вычисление оценки правдоподобности фактов, применение функции потерь (например, NLLLoss) и оптимизацию эмбеддингов. В результате обучения TComplEx генерирует обученные матрицы эмбеддингов для сущностей (U), отношений (R) и времени (W).

##### 2.3.1.4.1.2 Модель Question Answering (EmbedKGQA)

Модель EmbedKGQA использует эмбеддинги TComplEx для ответов на естественно-языковые вопросы. Она кодирует вопросы с помощью DistilBERT, проецирует их в "псевдо-отношения" (q_e_ent для сущностей и q_e_time для времени) в пространстве эмбеддингов графа знаний, оценивает кандидатов, объединяет оценки для ранжирования ответов и обучается с использованием NLLLoss.

#### 2.3.1.4.2 Анализ рабочего процесса

Рабочий процесс проекта детально представлен в блок-схеме `formal_workflow_v4.html`. Данная диаграмма иллюстрирует последовательность этапов от подготовки данных до расчета потерь обучения, охватывая как процесс обучения TComplEx, так и функционирование QA-модели.

На первом этапе, **"Формат Данных для Обучения"**, описывается структура входных данных, включающая исходный граф знаний, вопросы и данные, специально обработанные для TComplEx.

**"Этап 1: Обучение TComplEx"** детализирует процесс получения комплексных эмбеддингов для сущностей, отношений и времени. Этот этап включает генерацию негативных примеров, вычисление оценок правдоподобности фактов и оптимизацию эмбеддингов с использованием функции потерь.

После обучения TComplEx, начинается работа QA-модели. **"Этап 2.0: Подготовка данных (Предобработка)"** описывает токенизацию естественно-языковых вопросов с помощью `DistilBertTokenizer`.

На **"Этапе 2.1: Кодирование вопроса"** токенизированный вопрос преобразуется `DistilBERT` в векторное представление, которое затем проецируется в два псевдо-отношения, необходимые для взаимодействия с эмбеддингами графа знаний.

**"Этап 2.2: Оценка Кандидатов"** демонстрирует, как эти псевдо-отношения используются для оценки релевантности сущностей и временных меток в графе знаний.

На **"Этапе 2.3: Предсказание"** происходит объединение оценок кандидатов и их преобразование в ранжированный список ответов.

Завершающий **"Этап 2.4: Расчет Потери Обучения"** описывает вычисление функции потерь для обучения QA-модели, что позволяет корректировать параметры модели на основе расхождения между предсказанными и истинными ответами.

#### 2.3.1.4.3 Используемый набор данных

В проекте используется набор данных `wikidata_big`. Этот набор данных представляет собой крупномасштабный временной граф знаний, извлеченный из Wikidata. Он содержит информацию о сущностях, отношениях между ними и временных метках, указывающих на период истинности фактов. Набор данных разделен на обучающую, валидационную и тестовую выборки. Для удобства работы с моделью TComplEx, исходные текстовые данные преобразуются в числовые идентификаторы, а также создаются маппинги для обратного преобразования ID в текстовые метки.

#### 2.3.1.4.4 Метрики оценки и результаты

Оценка производительности модели проводилась с использованием стандартных метрик для задач KGC и QA: Hits@k и Mean Reciprocal Rank (MRR). Анализ результатов, представленных в `analasis.txt`, показывает следующее:

**Общие метрики на тестовой выборке:**
-   **Loss:** 1703.860312
-   **Hits@1:** 0.641
-   **Hits@10:** 0.875

**Детализация по типам вопросов (Hits@1):**
-   `simple_entity`: 0.988 (очень высокая точность для простых вопросов о сущностях)
-   `simple_time`: 0.985 (очень высокая точность для простых вопросов о времени)
-   `time_join`: 0.465
-   `first_last`: 0.371
-   `before_after`: 0.288

**Детализация по сложности вопросов (Hits@1):**
-   `simple`: 0.987
-   `complex`: 0.382

**Детализация по типу ответа (Hits@1):**
-   `entity`: 0.69
-   `time`: 0.549

**Детализация по типам вопросов (Hits@10):**
-   `simple_entity`: 0.993
-   `simple_time`: 0.991
-   `time_join`: 0.78
-   `first_last`: 0.814
-   `before_after`: 0.666

**Детализация по сложности вопросов (Hits@10):**
-   `simple`: 0.992
-   `complex`: 0.788

**Детализация по типу ответа (Hits@10):**
-   `entity`: 0.885
-   `time`: 0.857

**Интерпретация результатов:**
Модель демонстрирует очень высокую производительность на простых вопросах (`simple_entity`, `simple_time`), достигая почти 99% Hits@1. Однако, на более сложных типах вопросов, таких как `before_after`, `first_last` и `time_join`, а также на вопросах, классифицированных как `complex`, производительность значительно снижается. Это указывает на то, что модель успешно справляется с извлечением прямых фактов, но испытывает трудности с более сложными временными рассуждениями или многошаговыми запросами. Разница между Hits@1 и Hits@10 показывает, что правильный ответ часто находится в топ-10 предсказаний, но не всегда является первым.
\end{document}